\subsection{全四単純スペクトルの四曲線方向と有界有限部分への必要被覆}
\label{subsec:cubic-four-simple-spectral-curve-atlas}
\begin{lemma}[二次・三次の指数平面では四次非退化性を満たさない]
\label{lem:cubic-four-simple-extra-quartic}
$u$ を21二次共鳴の指数、$v\in\mathbb N^4,|v|=3$ とする。
$\mathbb Q u+\mathbb Q v$ に属する非負四次指数の集合は、四つの座標全てを被覆しない。
従ってH3の四単純スペクトルには、この二指数から独立な四次top指数が必要である。
\end{lemma}
\begin{proof}
座標置換と $u$ の非零scalar倍を除き、二次共鳴指数は
\[
 (1,-1,0,0),\quad (2,-1,-1,0),\quad (1,1,-1,-1)
\]
の三型である。
$w=au+bv,|w|=4$ なら $b=4/3$。
$u_j=0$ の座標があれば $w_j=4v_j/3$ なので、非負整数性から
$v_j=0$ 又は3である。前者なら全 $w$ がこの座標を欠く。
後者なら $v=3e_j$ で、$w=4e_j+au$ の他の座標の非負性は $a=0$ を強制し、やはり被覆しない。

最後の型では、$u_i=1,u_j=-1$ の全てのpairについて
$w_i+w_j=4(v_i+v_j)/3$ なので $v_i+v_j$ は3の倍数である。
$v$ が純三乗でなければ、この条件は不可能である。
実際、cross和3があれば他二座標は零で、正な1又は2と反対符号側の零とのcross和が条件に反する。
全cross和0なら $v=0$ となる。
従って $v=3e_j$ だけが残り、他三座標に正負の $u$ があるので再び $w=4e_j$ だけとなる。
四次topは全固有変数を含まなければならないから、独立な四次指数が必要である。
\end{proof}

\begin{lemma}[top重み零の一次元torusの二つのmoment条件]
\label{lem:cubic-torus-regrading-moments}
14641型のquadratic Koszul AS4整域 $B$ が、生成元の重み
$n_1,\ldots,n_4\in\mathbb Z$ のtorus作用を持ち、四次topの重みが零であるとする。
二次関係の重み $s$ の次元を $r_s$ とすれば
\[
 \sum_i n_i=0,\qquad r_s=r_{-s},\qquad
 \sum_s r_s s^2=2\sum_i n_i^2.
\]
\end{lemma}
\begin{proof}
middleのperfect pairingから $r_s=r_{-s}$、従って $\sum_s r_s s=0$。
$N>\max_i|n_i|$ を任意の整数とし、生成元の新次数を $N+n_i$ に取る。
全関係は斉次で、等変Koszul解消のEuler分母は
\[
 P_N(t)=1-\sum_i t^{N+n_i}+\sum_s r_st^{2N+s}
          -\sum_i t^{3N-n_i}+t^{4N}.
\]
旧次数 $d$ の各元の新次数は $(N-\max|n_i|)d$ と $(N+\max|n_i|)d$ の間にある。
従って $0<t<1$ で新Hilbert級数は、旧 $1/(1-t)^4$ にこの二つのdegree比を代入した級数で上下から挟まれる。
$1/P_N(t)$ は $t=1$ で正確に4次のpoleを持つ。
特に $(t\partial_t)^2P_N(1)=0$ であり、展開すると
\[
 4N\sum_i n_i+\sum_s r_s s^2-2\sum_i n_i^2=0.
\]
全ての大きな $N$ に対して成立するので二つの係数が零となる。
有限位数の中山自身がtorusを生成すると仮定した証明ではない。
\end{proof}

\begin{theorem}[四曲線方向と有界有限位数部分による全必要被覆]
\label{thm:cubic-four-simple-four-curve-cover}
強いintrinsic H3包絡の中山が四つの相異なる固有値を持つとする。
固有値のラベルを置換すれば、次の少なくとも一方が成立する。
\begin{enumerate}
\item ある $1\le N\le96$ に対して、全固有値が $N$ 乗根である。
\item ある $1\le M\le24$ と $\alpha_i^M=1$、$t\ne0$ によって
\[
 (\lambda_1,\ldots,\lambda_4)=(\alpha_1t^{n_1},\ldots,\alpha_4t^{n_4})
\]
と書け、方向 $n$ は次の四つのいずれかである：
\[
 (0,0,1,-1),\quad (0,1,-3,2),\quad
 (0,1,-2,1),\quad (1,-2,-1,2).
\]
\end{enumerate}
$M$ は24の約数とは限らず、$t$ は有限位数でもよい。
これは有限個の明示的なphase曲線と有界有限集合による必要被覆であり、存在・AS4十分性・H3許容性又は幾何学的群oidの分類ではない。
\end{theorem}
\begin{proof}
四単純の二次非共鳴部分は既分類skewとなり、三固有値以上を持たないので除外される。
従って二次共鳴を一つ選び、その指数を $u$、右辺を $\epsilon\in\{1,-1\}$ とする。
最初の型では $\epsilon=-1$、後二型では1である。
正常三次 $G$ の実際の非零語の全指数と、四次topの全非零語の指数を取り、
それらと $u$ の有理spanを $L$ とする。
前補題から $\dim L\ge3$ であり、三次指数 $v$ と四次指数 $w$ を選んで
$u,v,w$ を独立にできる。

$\dim L=4$ なら、さらに三次又は四次指数を一つ選び、整数 $4\times4$ 行列を可逆にできる。
その行列式の絶対値を $d$ とする。
monomial方程式へadjugateを適用すると、後二型では全 $\lambda_i^d=1$、最初の型では全 $\lambda_i^{2d}=1$。
全三次20指数、全四次35指数の有限列挙による最大値は
\[
\begin{array}{c|r|r|r}
 u&\max|\det(u,v,w,w')|&\max|\det(u,v,w,v')|&
 \max\gcd(3\text{次minors}(u,v,w))\\\hline
 (1,-1,0,0)&48&36&12\\
 (2,-1,-1,0)&96&72&24\\
 (1,1,-1,-1)&48&36&12
\end{array}
\]
である。従って第一項の $N\le96$ を得る。
列挙は整数3次minorのcofactor vectorと最後の行との内積だけで再生でき、浮動小数点を使わない。

$\dim L=3$ とする。$L$ のprimitive整数核を $n$ とする。
重要なのは、中山の有限・無限位数にかかわらず、実際の全top語と全 $G$ 語が $n$ 重み零となることである。
middle像が全二次関係なので、$B$ は実際にこのtorus作用を持ち、$G$ も固定される。
前補題から $\sum n_i=0$ とmoment条件を得る。

全 $21\cdot20\cdot35$ 個の $(u,v,w)$ のcofactor vectorをprimitive化し、
四次の重み零指数が全変数を被覆するものだけを残すと、符号正規化された方向は1062個。
$\sum n_i=0$ を課し、座標置換とparameter反転でまとめると次の十方向だけである。
以下 $r_{-s}=r_s$ で、表の三次は可換指数の略記であり、語を可換とした等式ではない。
\[
\begin{array}{c|c|l|l}
 n&r_0&(s:r_s),\ s>0&\text{重み零三次指数}\\\hline
 (0,0,1,-1)&2&1:2&x^3,x^2y,xy^2,y^3,xbc,ybc\\
 (0,1,-3,2)&0&1:1,2:1,3:1&x^3,ybc\\
 (0,1,-2,1)&0&1:2,2:1&x^3,y^2b,ybc,bc^2\\
 (1,-6,2,3)&0&3:1,4:1,5:1&yc^2\\
 (1,-5,-2,6)&0&1:1,4:1,7:1&x^2b\\
 (1,-4,-3,6)&0&2:1,3:1,7:1&b^2c\\
 (1,-4,1,2)&0&2:1,3:2&yc^2\\
 (1,-3,-2,4)&0&1:1,2:1,5:1&x^2b,b^2c\\
 (1,-2,-2,3)&0&1:2,4:1&x^2y,x^2b\\
 (1,-2,-1,2)&2&1:1,3:1&x^2y,b^2c
\end{array}
\]
関係数の列も必要条件から一意に得られる。
実際、torus固有空間の次元を $c_a$、二次tensorの重み $s$ の次元を $d_s$ とすると、
整域性とSegre多様体の射影交叉次元から
\[
 0\le r_s\le d_s-\max_{a+b=s}(c_a+c_b-1)
\]
である。右辺の各 $U_a\otimes U_b$ にrank-one tensorのSegre多様体があり、
これと交わる関係部分空間は零因子を与えるからである。
この上界、$\sum r_s=6$、対称性と二次momentを全整数値で解くと表の唯一の解になる。
公開検算器が全方向と全整数値を厳密に列挙する。

中ほどの四つの単一三次指数の方向は直ちに除外できる。
その二文字のpair重み空間は二次元で関係数1なので、整域性から両係数非零のskew関係がある。
従ってその三次の全語は一つの単項式に比例し、全 $G$ が一次因子を持つ。
$(1,-2,-2,3)$ では $x$ と $\langle y,b\rangle$ のmixed二次関係数が2で、
両射影の核は純積となるので可逆graphである。
$x,c$ もskew可換であり、$x$ は全包絡で正常となる。
全三次指数は $x^2\langle y,b\rangle$ なので、$G$ は共通一次因子 $x$ を持つ。

$(1,-3,-2,4)$ でも $x,b$ はskew可換。
重み2の関係は
\[
 a\,bc+b'\,cb+c'\,x^2=0
\]
で、$a,b'$ の少なくとも一方は非零である。
$a\ne0$ なら $b^2c$ を $bcb$ と $bx^2$ に書き換えられ、
全候補語は右因子 $b$ を持つ。
$a=0$ なら $b'\ne0$ であり、$cb$ を $x^2$ に書き換えれば全候補語は左因子 $b$ を持つ。
$x,b$ のskew可換性を使うと $x^2b$ も同じ因子を持つ。
従ってこれら六方向は全てH3整域商を与えず、表示した四方向だけが残る。

最後に $\lambda^u=\epsilon,\lambda^v=\lambda^w=1$ のrank3整数行列へSmith標準形を適用する。
各連結成分はprimitive $n$ の一次元torusのtranslateであり、torsion代表 $\alpha$ を選べる。
最大Smith因子は3次minorsのgcd以下なので、最初の型では代表の共通位数はそのgcdの2倍以下、
他二型ではgcd以下である。上の表から $M\le24$ を得る。
この議論は $t$ がroot of unityの点を除かない。
有限位数の中山自身からtorusを作る代わりに、全support指数のrank3性から実際のtorus作用を構成した。
以上で全必要被覆を得た。
\end{proof}

公開再現器は \texttt{verify\_cubic\_four\_simple\_spectral\_curve\_atlas.py}。
全60指数平面、全rank3方向、十方向の全関係rank、因子化に必要な孤立pair、
全整数determinantとSmith因子上界を再生する。
正のregradingによるpole、domain・因子、torus成分の通常証明は本文で与えた。
四曲線の全phaseと有限位数点の正常三次incidence、全CY・成分・境界・相対降下は未完成である。
